% Section 5.5, from lectures/L05/appendix/d_exercises.md.

\section{Exercises}
\label{c5:sec:exercises}\label{c5:app:d}

\begin{aside}
\textbf{How to work these.} These exercises practice the concepts from \secref{c5:sec:function} to
\secref{c5:sec:class}. Each set is one program: create a directory for it with the Makefile from
\secref{c1:sec:makefile} (or the one from \secref{c1:sec:tips}, for a program with an
\file{include} directory), and build and run it with \code{make}. Check your output against the
expected output each exercise gives.

The solutions are in \secref{sol:sec:c5}. Try each exercise yourself before looking at them.
\end{aside}

% ------------------------------------------------------------------------------------------
\exerciseset{Function Templates}
\label{c5:set:1}

\exercise{Clear bit}{Code}
\label{c5:ex:1.1}
Create a function template to clear a bit in a register:

\begin{cppcode}
/**
 * @brief Clear bit in a register.
 *
 * @tparam T Register type. Must be integral.
 *
 * @param[in, out] reg Destination register.
 * @param[in] bit Bit to clear.
 */
template<typename T>
constexpr void clear(T& reg, const std::uint8_t bit) noexcept;
\end{cppcode}

\task{Tasks}
\begin{parts}
  \item Use \code{static_assert()} with \code{std::is_integral<T>::value} and the following error
    message:
\begin{console}
Cannot perform bit operation with non-integral type!
\end{console}
  \item Clear the selected bit.
  \item Try the \code{clear()} function with the following code:
\end{parts}

\begin{cppcode}
#include <bitset>
#include <cstdint>
#include <iostream>

int main()
{
    std::uint8_t reg{0xFFU};

    // Clear bit 2 => reg becomes 0b11111011.
    clear(reg, 2U);
    std::cout << "Register content after clearing bit 2: 0b" << std::bitset<8U>(reg) << "\n";
    return 0;
}
\end{cppcode}

\noindent Expected output:

\begin{console}
Register content after clearing bit 2: 0b11111011
\end{console}

\exercise{Toggle bits}{Code}
\label{c5:ex:1.2}
Create the following function template to toggle one or several bits in a register:

\begin{cppcode}
/**
 * @brief Toggle bits in a register.
 *
 * @tparam T Register type. Must be integral.
 * @tparam Bits Parameter pack of bits.
 *
 * @param[in, out] reg Destination register.
 * @param[in] bits Bits to toggle.
 */
template<typename T, typename... Bits>
constexpr void toggle(T& reg, const Bits... bits) noexcept;
\end{cppcode}

\task{Tasks}
\begin{parts}
  \item Use a parameter pack.
  \item Ensure that \code{T} is an integral type using \code{static_assert()}.
  \item Iterate over the bits (for example using \code{{bits...}}).
  \item Toggle the specified bits.
  \item Try the \code{toggle()} function with the following code snippet:
\end{parts}

\begin{cppcode}
#include <bitset>
#include <cstdint>
#include <iostream>

int main()
{
    std::uint8_t reg{0xFFU};

    // Clear bit 2 => reg becomes 0b11111011.
    clear(reg, 2U);
    std::cout << "Register content after clearing bit 2: 0b" << std::bitset<8U>(reg) << "\n";

    // Toggle bits 0-3 => reg becomes 0b11110100.
    toggle(reg, 0U, 1U, 2U, 3U);
    std::cout << "Register content after toggling bits 0-3: 0b" << std::bitset<8U>(reg) << "\n";
    return 0;
}
\end{cppcode}

\noindent Expected output:

\begin{console}
Register content after clearing bit 2: 0b11111011
Register content after toggling bits 0-3: 0b11110100
\end{console}

% ------------------------------------------------------------------------------------------
\exerciseset{Function Template Concepts}
\label{c5:set:2}

\exercise[fillaccent2]{Template constraints}{Reflection}
\label{c5:ex:2.1}
\begin{parts}
  \item Explain why \code{static_assert()} is useful.
  \item What happens if the constraint fails?
  \item Why use compile-time checks instead of runtime?
\end{parts}

% ------------------------------------------------------------------------------------------
\exerciseset{Class Template Specialization}
\label{c5:set:3}

\exercise{Timer driver template}{Code}
\label{c5:ex:3.1}
In this exercise, you will create a timer driver that supports multiple implementations using
template specialization in a file \file{driver/timer/timer.hpp}.

The goal is to select the timer implementation at compile time based on a template parameter.

The timer type shall be selected using the following enumeration class:

\begin{cppcode}
namespace driver::timer
{
/**
 * @brief Enumeration of supported timers.
 */
enum class Type : std::uint8_t
{
    Stm32, ///< STM32 timer.
    Stub,  ///< Stub timer.
};
} // namespace driver::timer
\end{cppcode}

\noindent The timer shall be declared according to the following structure:

\begin{cppcode}
namespace driver::timer
{
/**
 * @brief Timer implementation.
 *
 * @tparam T Timer type.
 */
template<Type T>
class Timer;

} // namespace driver::timer
\end{cppcode}

\note \code{final} belongs on the definition of the class, not on this declaration. Written as
\code{class Timer final;}, the compiler would read \code{final} as the name of a variable of type
\code{Timer}.

The primary template shall represent a stub timer implementation. A specialization shall later be
created for \code{Type::Stm32}.

\task{Part I: Stub timer}
\begin{parts}
  \item Implement the primary template as a stub timer:
\begin{cppcode}
/**
 * @brief Timer implementation.
 *
 * @tparam T Timer type (default = stub timer).
 */
template<Type T = Type::Stub>
class Timer final
{
};
\end{cppcode}

  \item Add four private member variables:
  \begin{itemize}
    \item The first member variable shall:
    \begin{itemize}
      \item Be named \code{myTimeout_ms}.
      \item Represent the timeout duration in milliseconds.
      \item Have the type \code{std::uint16_t}.
      \item Be readable only after initialization.
    \end{itemize}
    \item The second member variable shall:
    \begin{itemize}
      \item Be named \code{myCounter_ms}.
      \item Represent the internal timer counter.
      \item Have the type \code{std::uint16_t}.
    \end{itemize}
    \item The third member variable shall:
    \begin{itemize}
      \item Be named \code{myRunning}.
      \item Represent whether the timer is active or not (\code{true/false}).
      \item Be \code{true} when the timer is running.
    \end{itemize}
    \item The fourth member variable shall:
    \begin{itemize}
      \item Be named \code{myInitialized}.
      \item Represent whether the timer is initialized or not (\code{true/false}).
      \item Be \code{true} when the timer is initialized.
    \end{itemize}
  \end{itemize}

  \item Add a constructor that:
  \begin{itemize}
    \item Takes a timeout value in milliseconds (\code{std::uint16_t}).
    \item Uses an initialization list to initialize the member variables.
    \item Initializes the counter to \code{0}.
    \item Initializes the timer as stopped.
    \item Marks the timer as initialized if the timeout value is valid (\code{> 0}).
    \item Is marked \code{explicit} and \code{noexcept}.
    \item If the timeout is valid, prints:
\begin{console}
Created stub timer with timeout 1000 ms!
\end{console}
      where \code{1000} shall be replaced by the configured timeout value.
    \item If the timeout is 0, prints:
\begin{console}
Failed to initialize stub timer: invalid timeout 0 ms!
\end{console}
  \end{itemize}

  \item Add a destructor that:
  \begin{itemize}
    \item Is marked \code{noexcept}.
    \item Stops the timer if it is currently running.
    \item Prints \code{Destroying stub timer!}.
  \end{itemize}

  \item Define the following methods:
  \begin{itemize}
    \item The method \code{timeout_ms()} shall:
    \begin{itemize}
      \item Return the configured timeout in milliseconds.
      \item Take no parameters.
      \item Be marked \code{const}, \code{noexcept}, and \code{[[nodiscard]]}.
    \end{itemize}
    \item The method \code{isRunning()} shall:
    \begin{itemize}
      \item Return \code{true} if the timer is running; otherwise \code{false}.
      \item Take no parameters.
      \item Be marked \code{const}, \code{noexcept}, and \code{[[nodiscard]]}.
    \end{itemize}
    \item The method \code{isInitialized()} shall:
    \begin{itemize}
      \item Return \code{true} if the timer is initialized; otherwise \code{false}.
      \item Take no parameters.
      \item Be marked \code{const}, \code{noexcept}, and \code{[[nodiscard]]}.
    \end{itemize}
    \item The method \code{start()} shall:
    \begin{itemize}
      \item Start the timer if the timer is initialized.
      \item Print \code{Starting stub timer!}.
      \item Take no parameters.
      \item Return no value.
      \item Be marked \code{noexcept}.
    \end{itemize}
    \item The method \code{stop()} shall:
    \begin{itemize}
      \item Stop the timer if the timer is initialized.
      \item Print \code{Stopping stub timer!}.
      \item Take no parameters.
      \item Return no value.
      \item Be marked \code{noexcept}.
    \end{itemize}
    \item The method \code{toggle()} shall:
    \begin{itemize}
      \item Toggle the timer if the timer is initialized.
      \item Print \code{Toggling stub timer!}.
      \item Take no parameters.
      \item Return no value.
      \item Be marked \code{noexcept}.
    \end{itemize}
    \item The method \code{tick()} shall:
    \begin{itemize}
      \item Simulate that 1 millisecond has passed.
      \item Increment the counter if the timer is running; otherwise do nothing.
      \item Return no value.
      \item Be marked \code{noexcept}.
    \end{itemize}
    \item The method \code{hasTimedOut()} shall:
    \begin{itemize}
      \item Return \code{true} when the counter has reached the timeout value; otherwise
        \code{false}.
      \item Reset the counter to \code{0} when a timeout occurs.
      \item Be marked \code{noexcept} and \code{[[nodiscard]]}.
    \end{itemize}
  \end{itemize}

  \item Delete the following methods:
  \begin{itemize}
    \item The default constructor.
    \item The copy constructor.
    \item The move constructor.
    \item The copy assignment operator.
    \item The move assignment operator.
  \end{itemize}
\end{parts}

\task{Part II: STM32 timer}
\begin{parts}
  \item Create a specialization for microcontroller \code{STM32}:
\begin{cppcode}
/**
 * @brief STM32 timer implementation.
 */
template<>
class Timer<Type::Stm32> final
{
};
\end{cppcode}

  \item Implement the STM32 timer with the same public interface as the primary template:
  \begin{itemize}
    \item Implement the same methods, but adjust the prints, for instance:
\begin{console}
Starting STM32 timer!
Stopping STM32 timer!
Destroying STM32 timer!
\end{console}
    \item Delete the default constructor as was done for the stub implementation.
    \item Prohibit copy and move semantics as was done for the stub implementation.
  \end{itemize}
\end{parts}

\task{Part III: Testing the timers}
Test the timers using the following test program in \file{main.cpp}:

\begin{cppcode}
#include <chrono>
#include <cstdint>
#include <thread>

#include "driver/timer/timer.hpp"

using namespace driver;

namespace
{
// -----------------------------------------------------------------------------
template<timer::Type T>
constexpr const char* timerType() noexcept
{
    if constexpr (T == timer::Type::Stub) { return "Stub"; }
    else if (T == timer::Type::Stm32) { return "STM32"; }
    return "Unknown";
}

// -----------------------------------------------------------------------------
template<timer::Type T>
void tickTimer(timer::Timer<T>& timer) noexcept
{
    constexpr const char* type{timerType<T>()};
    timer.tick();

    if (timer.hasTimedOut())
    {
        std::printf("%s timer has timed out after %u ms!\n", type, timer.timeout_ms());
    }
}

// -----------------------------------------------------------------------------
void delay_ms(const std::uint16_t ms) noexcept
{
    std::this_thread::sleep_for(std::chrono::milliseconds(ms));
}
} // namespace

// -----------------------------------------------------------------------------
int main()
{
    timer::Timer timer1{500U};
    timer::Timer<timer::Type::Stm32> timer2{1500U};
    timer1.start();
    timer2.start();

    constexpr std::uint16_t iterationCount{2000U};

    for (std::uint16_t i{}; i < iterationCount; ++i)
    {
        tickTimer(timer1);
        tickTimer(timer2);
        delay_ms(1U);
    }
    return 0;
}
\end{cppcode}

\noindent Expected output:

\begin{console}
Created stub timer with timeout 500 ms!
Created STM32 timer with timeout 1500 ms!
Starting stub timer!
Starting STM32 timer!
Stub timer has timed out after 500 ms!
Stub timer has timed out after 500 ms!
Stub timer has timed out after 500 ms!
STM32 timer has timed out after 1500 ms!
Stub timer has timed out after 500 ms!
Stopping STM32 timer!
Destroying STM32 timer!
Stopping stub timer!
Destroying stub timer!
\end{console}

\exercise[fillaccent2]{Reflection}{Reflection}
\label{c5:ex:3.2}
\begin{parts}
  \item Why is the timer type selected at compile time?
  \item What is a template specialization?
  \item What advantages does this design provide compared to using inheritance and virtual
    functions?
  \item When would inheritance be a better choice?
  \item What happens if the specialization for \code{Type::Stm32} is removed and a
    \code{Timer<Type::Stm32>} instance is created?
  \item Why might a template specialization be preferred over a virtual interface in a small
    microcontroller system?
\end{parts}
